\section*{第 \ref{ch:b3:sequence-stratigraphy} 章——层序地层学}

\begin{solution}{exo:b3:sequence-stratigraphy:1}
$A = 0.5 + 0.1 = \qty{0.6}{m/ka}$；在 \qty{20}{ka} 内，
可容纳空间为 \qty{12}{m}。海平面下降时，$A = -0.3 + 0.1 = \qty{-0.2}{m/ka}$，
表现为相对下降，但被沉降所减缓。
\end{solution}

\begin{solution}{exo:b3:sequence-stratigraphy:2}
(a) \omterm{def:b3:sequence-stratigraphy:tract}{下降期体系域}（\omterm{def:b3:sequence-stratigraphy:regression}{强迫海退}）。(b) \omterm{def:b3:sequence-stratigraphy:tract}{海侵体系域}。(c)
\omterm{def:b3:sequence-stratigraphy:tract}{高位体系域}。(d) \omterm{def:b3:sequence-stratigraphy:tract}{低位体系域}（下降之后的\omterm{def:b3:sequence-stratigraphy:regression}{正常海退}）。
\end{solution}

\begin{solution}{exo:b3:sequence-stratigraphy:3}
每个\omterm{def:b3:sequence-stratigraphy:parasequence}{准层序}的顶部都比其下方的\omterm{def:b3:sequence-stratigraphy:parasequence}{准层序}更浅，构成进积式叠置，
岸线向海推进，如高位或\omterm{def:b3:sequence-stratigraphy:tract}{低位体系域}中所见。与之相反的
顺序，即每个\omterm{def:b3:sequence-stratigraphy:parasequence}{准层序}都更细、更深，则为退积式叠置，
即\omterm{def:b3:sequence-stratigraphy:tract}{海侵体系域}的特征。
\end{solution}

\begin{solution}{exo:b3:sequence-stratigraphy:4}
每束含五个层对，即短偏心率（约
\qty{100}{ka}）与岁差（约 \qty{20}{ka}）之比：每个层对为一个岁差旋回，
每束为一个偏心率旋回。速率为
$1.6/100 = \qty{1.6}{cm/ka}$，厚 \qty{0.32}{m} 的层对约历时
\qty{20}{ka}，二者一致。
\end{solution}

\begin{solution}{exo:b3:sequence-stratigraphy:5}
$\mathrm dx_s/\mathrm dt = (4000 - 0.15\times20\,000)/50 = \qty{20}{m/ka}$,
向海方向。岸线静止，即 $AW = Q$ 时，$A = 4000/20\,000 =
\qty{0.2}{m/ka}$。
\end{solution}

\begin{solution}{exo:b3:sequence-stratigraphy:6}
$\mathrm dx_s/\mathrm dt = 4000/50 + 0.2/0.01 = 80 + 20 = \qty{100}{m/ka}$;
仅海平面下降一项就占 \qty{20}{m/ka}，即五分之一。
\end{solution}

\begin{solution}{exo:b3:sequence-stratigraphy:7}
岸线后退的条件为 $Q < AW = 0.4\times25\,000 = \qty{10000}{m^3/ka}$（每米）；
当 $Q = 6000$ 时，$(6000 - 10\,000)/10 = \qty{-400}{m/ka}$，相当于
每千年后退 \qty{0.4}{km}：临滨较低时，
海侵进展迅速。
\end{solution}

\begin{solution}{exo:b3:sequence-stratigraphy:8}
$1/94.9 - 1/123.9 = \num{0.00247}\text{ 循环}/\unit{ka}$：记录需长于
$1/0.00247 = \qty{406}{ka}$ 方能分开二者，取 \qty{600}{ka}（分辨率
0.00167）即足够。$1/94.9 - 1/98.9 = 0.00042$：所需长度约
\qty{2400}{ka}，是记录长度的四倍。
\end{solution}

\begin{solution}{exo:b3:sequence-stratigraphy:9}
$16.4/(20\times40.98) = \qty{0.0200}{m/ka}$，即 \qty{2.0}{cm/ka}；火山灰层
相隔 $3.1/0.020 = \qty{155}{ka}$。
\end{solution}

\begin{solution}{exo:b3:sequence-stratigraphy:10}
(a) $\tan\alpha = -0.4/200 = -0.002$：以 \ang{0.11} 下降。(b)
$0.2/50 = 0.004$：以 \ang{0.23} 上升。地震剖面跨度达数十
千米，垂直放大数十至数百倍，故该角度清晰可见；
长仅数百米的悬崖只显示一米或更小的抬升，
淹没于地层的变化之中。
\end{solution}

\begin{solution}{exo:b3:sequence-stratigraphy:11}
$1\,296\,000/54.725282 = \qty{23.68}{ka}$ ($g_5$); $/57.924383 =
\qty{22.37}{ka}$ ($g_2$); $/68.377912 = \qty{18.95}{ka}$ ($g_4$);
$/67.834313 = \qty{19.11}{ka}$（$g_3$）。$g_2 - g_5 = 3.199101''$ 每年，
\qty{405.1}{ka}；$g_4 - g_5 = 13.65263''$，\qty{94.9}{ka}。
\end{solution}

\begin{solution}{exo:b3:sequence-stratigraphy:12}
两种情形都在相同速率下产生 $0.3\times100 = \qty{30}{m}$ 的可容纳空间；供给相同时，
几何形态也相同，因为沉积物只响应相对海平面的变化。
全球海平面变化处处相同：其他边缘上
同时代的相同层序，指向第一种历史。
由盆地热历史与荷载历史构建的沉降史
（回剥法），或氧同位素等独立的全球海平面变化记录，
同样可将二者区分开来。
\end{solution}

\begin{solution}{pb:b3:sequence-stratigraphy:1}
\textbf{1.} $\dot E = 30\times2\pi/400\times\sin(2\pi t/400)$：最大为
\qty{0.471}{m/ka}，在 \qty{100}{ka} 和 \qty{500}{ka} 时上升，在
\qty{300}{ka} 时下降。

\textbf{2.} 转折端为 \qty{0.020}{m/ka}，距 \qty{20}{km} 处为 0.037，距
\qty{50}{km} 处为 0.062，距 \qty{120}{km} 处为 \qty{0.120}{m/ka}。

\textbf{3.} 当 $A < 0$，即 $\sin(2\pi t/400) < -0.0346/0.471 = -0.0734$ 时：$t =
200 + (400/2\pi)\arcsin0.0734 = \qty{204.7}{ka}$.

\textbf{4.} 取 $\sigma = 0.0617$，$t = 400 - 63.66\arcsin0.131 =
\qty{391.6}{ka}$.

\textbf{5.} 约 \qty{187}{ka}（对照 200）：全球海平面下降须先
超过沉降，并在低位期到来之前便不再占优。

\textbf{6.} $0.045\times400 = \qty{18}{m}$，与全球海平面变化幅度
\qty{60}{m} 相比，约为三分之一。

\textbf{7.} 自 50 至约 \qty{120}{ka} 向陆，向海至约
\qty{440}{ka}，向陆至约 \qty{560}{ka}，再向海至
\qty{600}{ka}（模式的转折点为 118.5、439.5 和
\qty{559}{ka}；每 \qty{50}{ka} 一档的表格只能给出其所在区间）。

\textbf{8.} $(6000 - 0.30\times15\,000)/40 = \qty{37.5}{m/ka}$，向海。

\textbf{9.} $6000/15\,000 = \qty{0.40}{m/ka}$.

\textbf{10.} $6000/35.7 + 0.424/0.012 = 168.1 + 35.3 = \qty{203}{m/ka}$,
与模式值之差在 \qty{0.5}{\percent} 以内。

\textbf{11.} 供给量为速度的 \qty{83}{\percent}。

\textbf{12.} $(-42.7 - 16.3)/(50.7 - 17.4) = \qty{-1.77}{m/km}$, \ang{0.10}.
自其时起，第一点已沉降 $0.0345\times395.5 =
\qty{13.6}{m}$，第二点沉降 $0.0623\times208.5 = \qty{13.0}{m}$：两者几乎
相同，因此就现今高程而言，下降量 \qty{59.0}{m} 与全球海平面
下降量仅相差 \qty{0.6}{m}。当时岸线处的
相对海平面下降量，即全球海平面变化减去下降期间的沉降，要更小一些。

\textbf{13.} 150 和 200：高位期；250、300、350：下降期；400：
低位期；450、500、550：海侵期（岸线在
440 之后不久后退）；600：高位期。

\textbf{14.} 204.7 和 \qty{391.6}{ka} 对应 204.5 和 391.5：两者相同，
在模式步长 \qty{0.5}{ka} 之内。

\textbf{15.} 在陆架上，是下降期间暴露的面：
自转折端至 \qty{17}{km} 处的高位期沉积顶面，延伸覆盖于
逐级下降的下降期前积层之上；在盆地中，对应的
整合面位于\omterm{def:b3:sequence-stratigraphy:tract}{下降期体系域}底部或顶部，此为约定。真实的
陆架会被河流切割成\omterm{def:b3:sequence-stratigraphy:condensed}{下切谷}，风化为土壤，
并部分遭到侵蚀。

\textbf{16.} $(18.5 - 52.8)/(559 - 439.5) = \qty{-0.29}{km/ka}$；
\omterm{def:b3:sequence-stratigraphy:surfaces}{最大海泛面}形成于其末端。

\textbf{17.} 岸线迅速掠过近乎平坦的沿海平原：
供给物被截留于其附近，分布于较宽的区域，
每处所得甚少。在前缘之外、约 \qty{55}{km} 以远的盆地方向，海底
处于饥饿状态：此处形成\omterm{def:b3:sequence-stratigraphy:condensed}{凝缩段}。

\textbf{18.} 分辨率 $1/30 = \num{0.033}\text{ 周期}/\unit{m}$；奈奎斯特频率 $1/(2\times
0.01) = \num{50}\text{ 周期}/\unit{m}$。

\textbf{19.} $7.5/2.0 = 3.75$、1 和 $0.81/2.0 = 0.405$；轨道比值为
$405/100 = 4.05$ 和 $41/100 = 0.41$。

\textbf{20.} \qty{2.0}{m} 对应 \qty{100}{ka}：\qty{2.0}{cm/ka}；$0.81/41 =
\qty{1.98}{cm/ka}$；$7.5/405 = \qty{1.85}{cm/ka}$。\omterm{def:b3:sequence-stratigraphy:periodogram}{周期图}在
\qty{8}{m} 附近能够显示的波长为 10、7.5 和 \qty{6}{m}：真实的
\qty{8.1}{m} 介于其中两者之间。

\textbf{21.} $30/0.020 = \qty{1500}{ka}$.

\textbf{22.} $23.7\times0.020 = \qty{0.474}{m}$。以
\qty{25}{cm} 为间隔采样时，奈奎斯特频率为每米 2 个周期，对应的波长为
\qty{0.5}{m}：岁差为每米 2.11 个周期，超出该频率，
折回为 $4 - 2.11 = 1.89$，即假波长 \qty{0.53}{m}。

\textbf{23.} 单个峰值总能与某一速率下的某一周期相吻合，无论其
成因为何；三个峰值的比值与轨道比值相符时，则与速率无关，
而由各峰值推出的速率必定相同。

\textbf{24.} 岸线由下降期开始时的 \qty{17.37}{km}
移至低位期海退结束时的 \qty{52.84}{km}，
即向盆地方向的最大位移为 \textbf{\qty[text-series-to-math=true]{35.5}{km}}。
\end{solution}
